\chapter{Introduction}

Pre-trained language models have changed how natural language processing systems are built. An encoder such as BERT~\cite{devlin2019bert} can be fine-tuned for a new classification problem with limited data, usually outperforming a task-specific network trained from scratch. An organisation supporting several language understanding tasks can therefore accumulate many fine-tuned copies of one architecture, each holding a complete parameter set even though most parameters stay close to the shared pre-trained values.

Multi-task learning~\cite{caruana1997multitask} offers the classical alternative of one model trained jointly on all tasks, but joint training requires access to all data at once and repeats its cost whenever a task is added. Keeping independent models avoids this but imposes the storage and maintenance burden described above. These two positions motivate \emph{model consolidation}: merging already fine-tuned models into one without returning to the full training data.

Simple Averaging and Task Arithmetic~\cite{ilharco2023task} merge parameters statically and therefore commit to one parameter vector for every input. Tang et al.~\cite{tang2024wemoe} addressed this limitation with Weight-Ensembling Mixture of Experts (WEMoE), in which a small router observes each input and determines how task-specific parameter updates are combined in the MLP blocks of a Transformer. Their work evaluated CLIP vision models on image classification and trained the routers through unsupervised test-time adaptation.

This report documents a project that adapts WEMoE to BERT and the GLUE benchmark and trains the routers with Few-Shot Fine-Tuning (FSFT) on a small labelled set instead of test-time adaptation.

The research objective is to determine whether input-conditioned MLP weights can recover task performance within one merged BERT model, and at what storage cost relative to five experts. Three questions follow. How much of the performance lost by static merging can routers recover when trained on labelled examples that the experts never saw? How does recovery depend on the labelled budget? And how does the choice of upscaled MLP modules affect accuracy, storage and per-task behaviour?

The contributions are practical, and no novel algorithm is claimed. The project (i) adapts the WEMoE router and weight-ensembling module to the post-norm BERT encoder with attention-mask-aware routing; (ii) applies a protocol that isolates expert training, router training, model selection and final reporting; (iii) evaluates two upscaling scopes and three labelled budgets; and (iv) analyses results against a majority-class floor, which exposes failures that an average conceals. The scope is limited to five binary GLUE tasks with one BERT-base backbone. Test-time adaptation is not implemented and is therefore out of scope.

Chapter~2 reviews the background, Chapter~3 presents the method as implemented, Chapter~4 specifies the experiments, Chapter~5 reports and discusses the results, and Chapter~6 concludes.